% Generated draft: requires mathematical and source review.
\documentclass[11pt]{article}
\usepackage[margin=25mm]{geometry}
\usepackage{fontspec}
\setmainfont{Latin Modern Roman}
\usepackage{amsmath,amssymb,mathtools,unicode-math}
\setmathfont{Latin Modern Math}
\usepackage{enumitem,needspace,xurl,hyperref}
\setlength{\parindent}{0pt}
\setlength{\parskip}{5pt}
\setlength{\emergencystretch}{3em}
\newcommand{\N}{\mathbb{N}}
\newcommand{\Z}{\mathbb{Z}}
\newcommand{\Q}{\mathbb{Q}}
\newcommand{\R}{\mathbb{R}}
\newcommand{\C}{\mathbb{C}}
\newcommand{\F}{\mathbb{F}}
\newcommand{\A}{\mathbb{A}}
\newcommand{\PP}{\mathbb{P}}
\newcommand{\E}{\mathbb{E}}
\newcommand{\eps}{\varepsilon}
\newcommand{\dd}{\,\mathrm{d}}
\begin{document}
\section*{UnsolvedMath ID: 30004425}
\textbf{Source problem code:} catalogue entry
\subsection{Definitions and mathematical statement}
Let \(T\) denote the tropicalization intended in the source, equipped with a polyhedral fan structure. Let \(C_+\) and \(C_-\) be distinct maximal-dimensional prime cones of \(T\) sharing a facet \(F\). Thus, writing \(d=\dim C_+=\dim C_-\), the adjacency hypothesis is
\[
C_+\cap C_-=F,\qquad \dim F=d-1.
\]
The ambient variety, ground field, embedding, and fan structure are not specified in the supplied record.

In an initial-ideal formulation, after choosing a field \(k\), an ambient polynomial or Laurent polynomial ring, a defining ideal \(I\), and a convention for initial forms, a cone is called prime when its associated initial ideal \(\operatorname{in}_w(I)\), constant for \(w\) in the relative interior of the cone, is prime. The supplied record does not specify these choices or the precise convention for a prime cone. A toric degeneration is a flat degeneration of the original variety with a toric special fiber; the source does not specify whether toric varieties must be normal or whether the construction is affine or projective.

An affine semigroup \(S\) is a finitely generated additive submonoid of a lattice, containing zero. Its semigroup algebra over \(k\) is
\[
k[S]=\bigoplus_{s\in S}k\chi^s,
\qquad \chi^s\chi^t=\chi^{s+t}.
\]
When an affine toric special fiber is presented by such an algebra, its semigroup can be described through \(S\). The supplied record does not identify the relevant semigroups, their lattices, or whether valuation semigroups or graded semigroups for projective constructions are intended.

\textbf{Problem as supplied.} Given the adjacent maximal-dimensional prime cones \(C_+\) and \(C_-\) and their associated toric degenerations, determine how the two degenerations are related across \(F\), including their wall-crossing maps and semigroups.

\textbf{Unresolved qualification.} This is a request to describe a relationship, rather than a specified conjectural identity or existence theorem. The construction of the associated degenerations, the domains and codomains of the wall-crossing maps, their required properties, and the intended relationship between the semigroups are not supplied. Consequently, the record does not determine a unique fully quantified mathematical assertion; the definitions above explain conventional terminology without supplying the missing hypotheses or selecting a particular wall-crossing construction.
\subsection{Short English statement}
The problem asks how the toric degenerations associated with two adjacent maximal-dimensional prime cones in a tropicalization are related across their common facet, particularly through wall-crossing maps and semigroups. The ambient setting and the maps and semigroups to be compared require specification. The supplied registry labels the problem open; this is a statement-only draft requiring human mathematical review.
\subsection{Sources}
\begin{itemize}
\item Local UnsolvedMath registry: supplied record for permanent ID \(30004425\), titled \emph{Wall Crossing Between Adjacent Tropical Toric Degenerations}; exact source problem code \texttt{catalogue entry}. The record belongs to the Oberwolfach Reports Open Problems collection and supplies proposed year \(2020\), without naming a proposer.
\item Supplied source citation: \emph{Mini-Workshop: Superpotentials in Algebra and Geometry} (2021), Oberwolfach Reports, pp.\ 569--600, DOI \texttt{10.4171/owr/2020/10}. Supplied source URL and literature-source link labeled ``original OWR report'': \url{https://doi.org/10.4171/owr/2020/10}. The collection describes these reports as documenting mathematical workshops and their open-problem discussions. No problem-specific page or named author is supplied.
\item Supplied literature assessment, checked on \texttt{2026-08-22}: status ``open''; the relationship across the facet is described as a natural question for which no general answer was identified in the primary literature checked. The listed source checked is the original OWR report linked above. This is supplied dated curation and literature triage, not an independently verified assessment; time-sensitive status and cited claims require independent verification.
\end{itemize}
\end{document}
